\documentclass[12pt]{article}
\usepackage[spanish]{babel}
\usepackage[utf8]{inputenc}
\usepackage[T1]{fontenc}
\usepackage{geometry}
\usepackage{setspace}
\usepackage{graphicx}
\usepackage[colorlinks=true, linkcolor=blue, citecolor=blue, urlcolor=blue]{hyperref}
\usepackage{listings}
\usepackage{xcolor}
\usepackage{float}
\usepackage{booktabs}
\usepackage{enumitem}
\definecolor{codebg}{RGB}{245,245,245}
\definecolor{codekeyword}{RGB}{0,0,180}
\definecolor{codecomment}{RGB}{95,140,95}
\definecolor{codestring}{RGB}{163,21,21}
\lstdefinestyle{bashStyle}{
    backgroundcolor=\color{codebg},
    commentstyle=\color{codecomment}\itshape,
    keywordstyle=\color{codekeyword}\bfseries,
    stringstyle=\color{codestring},
    basicstyle=\ttfamily\small,
    breakatwhitespace=false,
    breaklines=true,
    captionpos=b,
    keepspaces=true,
    numbers=none,
    showspaces=false,
    showstringspaces=false,
    showtabs=false,
    tabsize=2,
    frame=single,
    rulecolor=\color{gray!40},
    language=bash
}
\lstdefinestyle{sqlStyle}{
    backgroundcolor=\color{codebg},
    commentstyle=\color{codecomment}\itshape,
    keywordstyle=\color{codekeyword}\bfseries,
    stringstyle=\color{codestring},
    basicstyle=\ttfamily\small,
    breakatwhitespace=false,
    breaklines=true,
    captionpos=b,
    keepspaces=true,
    numbers=none,
    showspaces=false,
    showstringspaces=false,
    showtabs=false,
    tabsize=2,
    frame=single,
    rulecolor=\color{gray!40},
    language=SQL,
    morekeywords={SERIAL}
}
\lstset{style=sqlStyle}
\lstdefinestyle{panel}{
  backgroundcolor=\color{gray!8},
  basicstyle=\ttfamily\footnotesize,
  frame=single,
  framerule=0.4pt,
  rulecolor=\color{gray!50},
  numbers=none,
  breaklines=true,
  columns=fullflexible,
  showstringspaces=false,
  xleftmargin=12pt,
  xrightmargin=12pt,
  aboveskip=10pt,
  belowskip=10pt
}
\geometry{letterpaper, margin=1in}
\setstretch{1.25}
\usepackage[
  backend=biber,
  style=ieee,
  doi=true,
  url=true
]{biblatex}
\DeclareFieldFormat{doi}{\href{https://doi.org/#1}{#1}}
\addbibresource{Referencias.bib}
\setlength{\parindent}{0pt}
\setlength{\parskip}{6pt}
\begin{document}
\pagestyle{empty}
\renewcommand{\tablename}{Tabla}

% ============================================================
%  CARÁTULA
% ============================================================
\begin{titlepage}
\thispagestyle{empty}
\centering
\vspace*{0.2cm}

\includegraphics[width=3.2cm]{logo_uteq.png}

\vspace{0.7cm}

{\Large \textbf{UNIVERSIDAD TÉCNICA ESTATAL DE QUEVEDO}}\par
{\small TEMA:}\par
\vspace{0.1cm}
{\large \textbf{Tercera Entrega del Proyecto Fin de Curso (PFC) \\ Sistema de Gestión de Recursos Operativos, Administrativos y de Seguridad (SGROAS)}}\par

\vspace{0.2cm}
{\footnotesize GRUPO CET:}\par
\vspace{0.2cm}
{\small \textbf{CASTRO ESPINOZA KEVIN MOISES}}\par
{\small \textbf{ESCUDERO PLAZA MARIA DEL ROSARIO}}\par
{\small \textbf{TEJADA BAJAÑA LUIS ALEJANDRO}}\par

\vspace{0.2cm}
{\small \textbf{DOCENTE:}}\par
\vspace{0.2cm}
{\small \textbf{DR. GLEISTON CICERÓN GUERRERO ULLOA, PH.D.}}\par

\vspace{0.2cm}
{\small \textbf{CURSO:}}\par
\vspace{0.2cm}
{\small \textbf{5TO SOFTWARE A}}\par

\vspace{0.2cm}
{\small \textbf{FECHA:}}\par
\vspace{0.2cm}
{\small \textbf{VIERNES, 24 DE JULIO DEL 2026}}\par

\vspace{0.2cm}
{\small \textbf{LINK GITHUB:}}\par
\vspace{0.2cm}
{\small \texttt{https://github.com/Alxjandr07/SGROAS-ProyectoAppWeb.git}}\par

\vfill
\end{titlepage}

\pagestyle{plain}
\setcounter{page}{1}

% ============================================================
%  TABLA DE CONTENIDOS
% ============================================================
\tableofcontents
\clearpage

% ============================================================
%  RESUMEN EJECUTIVO
% ============================================================
\section*{Resumen ejecutivo}
\addcontentsline{toc}{section}{Resumen ejecutivo}

El presente informe documenta el estado del sistema SGROAS (Sistema de Gestión de Recursos
Operativos, Administrativos y de Seguridad) en su hito de Tercera Entrega, correspondiente a la
etiqueta \texttt{v0.9.0-rc} del repositorio \cite{sgroas_repo}. El documento sigue la estructura y
los requisitos técnicos establecidos en la guía institucional de la Tercera Entrega del Proyecto
Fin de Curso (PFC) \cite{uteq2026guia}, consolidando evidencia empírica reproducible sobre el
arranque del sistema, la documentación de la API, el flujo de autenticación JWT, el rendimiento
bajo carga, la cobertura de pruebas y el versionado semántico del repositorio.

% ============================================================
\section{Estado del sistema}

El sistema SGROAS se encuentra operativo en su totalidad mediante orquestación con Docker
Compose. La Figura~\ref{fig:imagen1} evidencia el arranque exitoso de los tres servicios que
componen el sistema (backend, PostgreSQL y Redis) desde un entorno Windows con Docker Desktop.

\begin{figure}[H]
    \centering
    \includegraphics[width=0.85\linewidth]{imagen1}
    \caption{Full system startup through \texttt{docker compose up}: backend image build and
    startup of the \texttt{sgroas-postgres}, \texttt{sgroas-redis-compose} and
    \texttt{sgroas-backend} containers.}
    \label{fig:imagen1}
\end{figure}

El estado de los contenedores en ejecución, con sus puertos publicados y tiempos de actividad,
se confirma en la Figura~\ref{fig:imagen8}, evidenciando que los tres servicios permanecen estables
durante la sesión de pruebas.

\begin{figure}[H]
    \centering
    \includegraphics[width=0.85\linewidth]{imagen8}
    \caption{\texttt{docker ps} output confirming the \texttt{sgroas-backend},
    \texttt{sgroas-postgres} and \texttt{sgroas-redis-compose} containers in \textit{Up}
    state, with the corresponding port mapping.}
    \label{fig:imagen8}
\end{figure}

La documentación interactiva de la API, generada automáticamente mediante Springdoc OpenAPI
3.1 y expuesta en \texttt{/api/docs}, se muestra en la Figura~\ref{fig:imagen2}, cumpliendo el
requisito del Bloque A.1 de la guía de entrega \cite{uteq2026guia}.

\begin{figure}[H]
    \centering
    \includegraphics[width=0.85\linewidth]{imagen2}
    \caption{SGROAS API Swagger UI (version \texttt{0.9.0-rc}, OAS 3.1 specification),
    showing the endpoints of the driver controller and the authentication controller.}
    \label{fig:imagen2}
\end{figure}

\subsection{Frontend Angular}

La interfaz de usuario del sistema se desarrolló con Angular 20 y se comunica con el backend
mediante cookies HttpOnly para la autenticación. La Figura~\ref{fig:login-frontend} muestra la
pantalla de inicio de sesión.

\begin{figure}[H]
    \centering
    \includegraphics[width=0.7\linewidth]{imagen-login}
    \caption{Angular frontend login screen with a reactive form and real-time validation.}
    \label{fig:login-frontend}
\end{figure}

El panel principal (dashboard) presenta los módulos del sistema con acceso directo a cada
sección, como se observa en la Figura~\ref{fig:dashboard}.

\begin{figure}[H]
    \centering
    \includegraphics[width=0.85\linewidth]{imagen-dashboard}
    \caption{Main SGROAS dashboard with an indicators summary and access to the six system
    modules.}
    \label{fig:dashboard}
\end{figure}

Las vistas de listado permiten la gestión paginada de conductores y usuarios, con opciones de
creación, edición y desactivación (Figura~\ref{fig:listado-conductores} y
Figura~\ref{fig:listado-usuarios}).

\begin{figure}[H]
    \centering
    \includegraphics[width=0.85\linewidth]{imagen-conductores}
    \caption{Paginated driver listing with edit and deactivation options.}
    \label{fig:listado-conductores}
\end{figure}

\begin{figure}[H]
    \centering
    \includegraphics[width=0.85\linewidth]{imagen-usuarios}
    \caption{Paginated system user listing showing role and status.}
    \label{fig:listado-usuarios}
\end{figure}

% ============================================================
\section{Arquitectura actual}

El sistema SGROAS sigue una arquitectura de tres capas: frontend Angular, backend Spring Boot
y base de datos PostgreSQL, con Redis como caché distribuida y Docker Compose para la
orquestación. Los diagramas C4 en niveles 1, 2 y 3 se encuentran versionados en
\texttt{docs/arquitectura/} del repositorio \cite{sgroas_repo}.

\subsection{Nivel 1 — Contexto del sistema}
El sistema interactúa con tres tipos de actores: Administrador, Coordinador y Seguridad,
cada uno con diferentes permisos de acceso.

\subsection{Nivel 2 — Contenedores}
Los contenedores principales son: aplicación Angular (frontend), API Spring Boot (backend),
base de datos PostgreSQL, caché Redis y el proxy de la aplicación.

\subsection{Nivel 3 — Componentes}
El backend se compone de controladores REST, servicios, repositorios JPA y la capa de
seguridad con JWT.

% ============================================================
\section{Reproducibilidad del artefacto}

Conforme al Bloque B de la guía de entrega \cite{uteq2026guia}, el sistema fija cada imagen de
contenedor por \textit{digest} \texttt{sha256} en lugar de una etiqueta variable, evitando
derivas silenciosas entre reconstrucciones. La configuración resuelta de Docker Compose se
presenta en la Figura~\ref{fig:imagen13}.

\begin{figure}[H]
    \centering
    \includegraphics[width=0.7\linewidth]{imagen13}
    \caption{\texttt{docker compose config} output showing the PostgreSQL and Redis image
    pinning by \textit{digest} \texttt{sha256}, together with the backend environment
    variables (JWT, database credentials and Redis).}
    \label{fig:imagen13}
\end{figure}

La correspondencia entre los \textit{digests} declarados y las imágenes descargadas localmente
se verifica en la Figura~\ref{fig:imagen14}.

\begin{figure}[H]
    \centering
    \includegraphics[width=0.85\linewidth]{imagen14}
    \caption{Local images listing with \texttt{docker images --digests}, filtered by
    \texttt{postgres} and \texttt{redis}, confirming the match of the \texttt{sha256}
    \textit{digest} used in \texttt{docker-compose.yml}.}
    \label{fig:imagen14}
\end{figure}

El versionado semántico estricto exigido por el Bloque E.4 \cite{uteq2026guia} se evidencia en
el historial de \textit{commits} y etiquetas del repositorio (Figura~\ref{fig:imagen11}), así como
en las \textit{releases} publicadas en GitHub (Figura~\ref{fig:imagen12}).

\begin{figure}[H]
    \centering
    \includegraphics[width=0.85\linewidth]{imagen11}
    \caption{\texttt{git log --oneline --decorate} history showing the \texttt{HEAD}
    \textit{commit} tagged as \texttt{v0.9.0-rc} on the \texttt{main} branch, preceded by
    the merge of the \texttt{feat/alejandro-auth} branch.}
    \label{fig:imagen11}
\end{figure}

\begin{figure}[H]
    \centering
    \includegraphics[width=0.85\linewidth]{imagen12}
    \caption{Tags published in the \texttt{Alxjandr07/SGROAS-ProyectoAppWeb} repository:
    \texttt{v0.9.0-rc} (Third Delivery) and \texttt{v0.7.1} (closing the Delivery 1B
    observations), the latter verified by GitHub.}
    \label{fig:imagen12}
\end{figure}

% ============================================================
\section{Matriz de trazabilidad}

La matriz de trazabilidad end-to-end (archivo \texttt{docs/trazabilidad/matriz.csv}) cubre el
100\,\% de los requisitos con prioridad \texttt{Must} vigentes en el estado \texttt{v0.9.0-rc}.
Cada requisito se traza desde su identificador hasta el módulo, endpoint, prueba automatizada,
tipo de acceso y evidencia empírica asociada. La Tabla~\ref{tab:trazabilidad} presenta un
resumen de los requisitos funcionales y no funcionales implementados.

\begin{table}[H]
\centering
\caption{Resumen de requisitos trazados en la matriz}
\label{tab:trazabilidad}
\begin{tabular}{@{}lcll@{}}
\toprule
\textbf{Tipo} & \textbf{Cantidad} & \textbf{Prioridad} & \textbf{Estado} \\
\midrule
Funcionales & 14 & Must & Implementado \\
No funcionales & 5 & Must & Verificado \\
No funcionales & 1 & Should & Implementado \\
\bottomrule
\end{tabular}
\end{table}

% ============================================================
\section{Protocolo experimental}

La evidencia empírica cuantitativa se organiza según el Bloque C de la guía de entrega
\cite{uteq2026guia}, cubriendo rendimiento, seguridad, usabilidad, cobertura de pruebas y
accesibilidad mediante corridas reproducibles.

\subsection{Flujo de autenticación evaluado}

El protocolo de prueba de autenticación se ejecutó desde Postman contra el endpoint
\texttt{POST /api/auth/login}, cubriendo tres escenarios: solicitud válida, solicitud con
error de validación y verificación del contenido del \textit{token} emitido.

% ============================================================
\section{Resultados por bloque}

\subsection{Rendimiento (Bloque C.1)}

Se ejecutaron tres corridas de carga con \textbf{k6} \cite{k6} contra el endpoint de listado
protegido, con 50 usuarios virtuales (VUs) durante 30 segundos cada una, conforme al protocolo
declarado en el Bloque B.2 y C.1 de la guía \cite{uteq2026guia}. Los resultados se presentan en
la Figura~\ref{fig:imagen7}.

\begin{figure}[H]
    \centering
    \includegraphics[width=0.85\linewidth]{imagen7}
    \caption{Result of the three k6 load runs (50 VUs, 30 s) against the listing endpoint:
    4468 HTTP requests in total, 0{,}00\,\% error rate, mean p95 response time 81{,}43\,ms and
    mean p99 188{,}16\,ms, meeting the \texttt{response\_time < 500 ms} threshold.}
    \label{fig:imagen7}
\end{figure}

Los dos \textit{checks} configurados (\texttt{status es 200} y \texttt{tiempo respuesta <
500ms}) se cumplieron en el 100\,\% de las 8930 verificaciones ejecutadas (4465 por \textit{check}
en cada una de las tres corridas), sin peticiones
fallidas (\texttt{http\_req\_failed} = 0{,}00\,\%). El reporte agregado con media,
desviación típica e intervalo de confianza al 95\,\% se encuentra en
\texttt{docs/mediciones/perf/REPORT.md}.

\subsection{Seguridad (Bloque C.2)}

Se realizó una auditoría automática de los seis controles OWASP mínimos usando scripts curl
disponibles en \texttt{docs/mediciones/sec/} \cite{owasp2021}.

\textbf{A01 — Control de acceso roto:} al intentar crear un conductor con un usuario de rol
SEGURIDAD (sin permisos), el sistema responde con \texttt{403 Forbidden}.

\textbf{A02 — Fallo criptográfico:} el servidor utiliza TLSv1.3 con cifrado AEAD
(\texttt{TLS\_AES\_256\_GCM\_SHA384}).

\textbf{A03 — Inyección:} el envío de payload \texttt{' OR '1'='1} en campos de búsqueda
responde con \texttt{422 Unprocessable Entity} y estructura ProblemDetails.

\textbf{A05 — Malconfiguración de seguridad:} las cabeceras HTTP incluyen
\texttt{Strict-Transport-Security}, \texttt{X-Frame-Options: DENY},
\texttt{X-Content-Type-Options: nosniff} y \texttt{Content-Security-Policy}.

\textbf{A07 — Fallo de identificación:} seis intentos fallidos consecutivos de login desde la
misma IP provocan una respuesta \texttt{429 Too Many Requests}.

\textbf{A09 — Registro y monitoreo:} los logs del backend registran cada intento de login con
timestamp, IP y email del usuario.

El módulo de autenticación se evaluó mediante peticiones directas con Postman contra
\texttt{POST /api/auth/login}, verificando el cumplimiento del estándar \textbf{JWT (RFC
7519)} \cite{rfc7519} y del formato de error \textbf{ProblemDetails (RFC 7807)}
\cite{rfc7807}.

\begin{figure}[H]
    \centering
    \includegraphics[width=0.85\linewidth]{imagen3}
    \caption{\texttt{POST /api/auth/login} request with valid credentials
    (\texttt{admin@sgroas.com}), returning \texttt{200 OK} in 1{,}83\,s and setting the
    \texttt{access\_token} cookie with the \texttt{HttpOnly} and \texttt{Secure} attributes
    enabled.}
    \label{fig:imagen3}
\end{figure}

\begin{figure}[H]
    \centering
    \includegraphics[width=0.85\linewidth]{imagen4}
    \caption{Request with an incomplete body (without the \texttt{password} field),
    returning \texttt{422 Unprocessable Entity} with a ProblemDetails structure as per RFC
    7807 \cite{rfc7807}: \texttt{type}, \texttt{title}, \texttt{status}, \texttt{detail} and
    \texttt{instance}, along with the per-field validation error detail.}
    \label{fig:imagen4}
\end{figure}

\begin{figure}[H]
    \centering
    \includegraphics[width=0.85\linewidth]{imagen5}
    \caption{\texttt{200 OK} response of the successful login, including \texttt{accessToken},
    \texttt{refreshToken}, \texttt{tokenType}, \texttt{expiresIn} and the authenticated user
    data (\texttt{name}, \texttt{email}, \texttt{role}).}
    \label{fig:imagen5}
\end{figure}

La estructura interna del JWT emitido se decodificó con la herramienta \textit{jwt.io} para
verificar la presencia de los siete \textit{claims} estándar exigidos por el Bloque A.1 de la
guía \cite{uteq2026guia,rfc7519} (Figura~\ref{fig:imagen6}).

\begin{figure}[H]
    \centering
    \includegraphics[width=0.85\linewidth]{imagen6}
    \caption{Decoding of the JWT issued by SGROAS showing the \texttt{jti}, \texttt{iss},
    \texttt{sub}, \texttt{aud}, \texttt{iat}, \texttt{nbf}, \texttt{exp} \textit{claims},
    together with the custom \texttt{name} and \texttt{role} \textit{claims}.}
    \label{fig:imagen6}
\end{figure}

\subsection{Usabilidad (Bloque C.3)}

Se aplicó el instrumento System Usability Scale (SUS) de Brooke \cite{brooke1996} mediante un
formulario electrónico a diez participantes externos al equipo, en sesiones moderadas en las que
cada participante operó la aplicación y respondió el cuestionario por sí mismo. Las tareas
evaluadas incluyen login, alta de un conductor, edición, eliminación lógica y logout. La
puntuación SUS media fue de \textbf{63{,}0} sobre 100 (DT = 13{,}9, IC~95\,\% = [53{,}1,
72{,}9]), lo que corresponde a la calificación adjetiva \textit{Bueno} en la escala de Bangor
et al.~\cite{bangor2009} y se ubica en la zona de aceptabilidad \textit{marginal} (50--70),
por debajo del umbral de 70 que dicha escala asocia a sistemas aceptables; este hallazgo se
documenta como área de mejora. La matriz de datos crudos se encuentra en
\texttt{docs/mediciones/sus/sus-raw.csv} y el desglose por participante en
\texttt{docs/mediciones/sus/REPORT.md}.

\subsection{Cobertura de pruebas (Bloque C.4)}

La cobertura de código se midió con \textbf{JaCoCo} \cite{jacoco} sobre las 50 clases de los
paquetes del dominio, DTOs, servicios, seguridad, controladores y configuración de la API,
mediante 150 pruebas unitarias de JUnit 5 (0 fallos, 0 errores). El reporte agregado se muestra
en la Figura~\ref{fig:imagen9} y la estructura de archivos generados en el repositorio, en la
Figura~\ref{fig:imagen10}.

\begin{figure}[H]
    \centering
    \includegraphics[width=0.85\linewidth]{imagen9}
    \caption{JaCoCo per-package report: 98{,}8\,\% global instruction coverage (3484 of 3527
    covered instructions), 85{,}4\,\% branch coverage (70 of 82) and 99{,}7\,\% line coverage
    (778 of 780). The \texttt{controller}, \texttt{dto}, \texttt{entity}, \texttt{config} and
    \texttt{exception} packages reach 100\,\% instruction coverage.}
    \label{fig:imagen9}
\end{figure}

\begin{figure}[H]
    \centering
    \includegraphics[width=0.5\linewidth]{imagen10}
    \caption{Versioned structure of the \texttt{docs/mediciones/jacoco/} directory in the
    repository, with the per-package breakdown and the \texttt{index.html}, \texttt{jacoco.csv}
    and \texttt{jacoco.xml} files required for evidence traceability.}
    \label{fig:imagen10}
\end{figure}

\noindent El plugin JaCoCo está configurado con una regla \texttt{check} que impone un mínimo de
60\,\% de cobertura tanto de líneas como de ramas (\texttt{BUNDLE}); la corrida de verificación
completa (150 tests) supera dicho umbral y el mensaje \texttt{All coverage checks have been met}
se obtiene en cada \texttt{mvn verify}, por lo que el criterio C4 de la guía \cite{uteq2026guia}
se cumple.

\subsection{Accesibilidad y calidad web (Bloque C.5)}

Se ejecutó Lighthouse CI contra el frontend Angular servido con compresión gzip
(configuración equivalente a producción, perfil móvil con \textit{throttling} de red
\textit{Slow 4G}). Los puntajes obtenidos fueron: Performance 100, Accessibility 95,
Best Practices 100 y SEO 90, cumpliendo con los umbrales mínimos establecidos en
\texttt{lighthouserc.js} (Performance \(\ge\) 80, Accessibility \(\ge\) 90, Best
Practices \(\ge\) 90, SEO \(\ge\) 90). Los archivos JSON crudos de cada corrida se
archivan en \texttt{docs/mediciones/lighthouse/lhci-*.json}.

% ============================================================
\section{Amenazas a la validez}

\begin{itemize}
    \item \textbf{Validez interna:} las tres corridas de k6 se ejecutaron en el entorno
    local de desarrollo, con media, desviación típica e intervalo de confianza al 95\,\%
    calculados sobre las corridas independientes \cite{uteq2026guia}; el entorno no refleja
    exactamente las condiciones de producción.
    \item \textbf{Validez de constructo:} la cobertura de JaCoCo (98{,}8\,\% de instrucciones,
    85{,}4\,\% de ramas y 99{,}7\,\% de líneas, con 150 pruebas JUnit 5) refleja el estado
    completo de las pruebas unitarias e de integración sobre los paquetes del dominio, servicios,
    seguridad, controladores y DTOs.
    \item \textbf{Validez externa:} las pruebas de seguridad (Figuras~\ref{fig:imagen3}
    a~\ref{fig:imagen6}) cubren el flujo de autenticación y los seis controles OWASP mínimos del
    Bloque C.2 \cite{owasp2021,uteq2026guia}; la evidencia cruda se encuentra en
    \texttt{docs/mediciones/sec/}.
    \item \textbf{Usabilidad:} la muestra de diez participantes es suficiente para estimaciones
    estables de la media SUS, pero no permite análisis segmentados por perfil de usuario; la media
    (63{,}0) se encuentra en la zona de aceptabilidad marginal, lo que motiva acciones de mejora
    de usabilidad para la entrega final.
\end{itemize}

% ============================================================
\section*{Declaración de contribuciones (CRediT)}
\addcontentsline{toc}{section}{Declaración de contribuciones (CRediT)}

Las contribuciones del equipo se documentan en el archivo \texttt{CONTRIBUTORS.md} del
repositorio \cite{sgroas_repo}, siguiendo la taxonomía CRediT de 14 roles conforme al
Bloque E.1 de la guía de entrega \cite{uteq2026guia}. El equipo autor de este informe está
conformado por Castro Espinoza Kevin Moisés, Escudero Plaza María del Rosario y Tejada Bajaña
Luis Alejandro.

% ============================================================
\section*{Declaración ética}
\addcontentsline{toc}{section}{Declaración ética}

El proyecto no procesa datos personales, financieros o de salud reales de terceros; los datos
semilla de la base de datos son sintéticos. La declaración completa se encuentra en
\texttt{docs/etica/ETHICS.md} conforme al Bloque F de la guía \cite{uteq2026guia}.

% ============================================================
\printbibliography[title={Referencias}]

\end{document}
